\section{Related Work}
\label{sec:related}

\subsection{Execution Feedback for Iterative Code Repair}

Iterative self-repair using execution feedback---generate code, execute it, use error output as context, re-generate---is a well-established approach. \citet{Chen2023Teaching} showed that execution trace feedback enables LLMs to self-debug, improving MBPP by 12\% and achieving 10$\times$ sample efficiency compared to best-of-N sampling. \citet{Shinn2023Reflexion} demonstrated that verbal reinforcement via execution results reaches 91\% HumanEval pass@1 in Reflexion. \citet{Gehring2024RLEF} showed that RL-grounded execution feedback (RLEF) further improves over prompting-only repair, reducing required samples by 10$\times$. Most recently, \citet{Arimbur2026HowMany} showed that modern 8B instruction-tuned models---including Llama 3.1 8B---achieve meaningful self-repair (+4.9 to +17.1pp HumanEval) with execution feedback alone, without fine-tuning, and that most gains occur in repair rounds 1--2.

These works collectively establish that execution feedback \emph{works}. However, none compare execution feedback against pylint/mypy static analysis on the same benchmarks under compute-controlled conditions. Our work fills this gap.

\subsection{Static Analysis Feedback for LLM Code Quality}

Pylint, mypy, and related static analysis tools have been integrated into LLM code generation pipelines primarily for quality metrics beyond functional correctness. \citet{Blyth2025Static} demonstrated that iterative pylint/bandit feedback reduces security issues in LLM-generated code from 40\% to 13\% over 10 iterations on PythonSecurityEval. This result establishes that pylint \emph{can} provide useful feedback---but on a benchmark where pylint's Error and Warning categories (security violations) are the dominant failure mode. HumanEval and MBPP failures are dominated by logic and runtime errors, not security issues, making this result difficult to generalize.

FeedbackEval \cite{Dai2025FeedbackEval} provides the most systematic comparison of feedback types, covering compiler feedback, test feedback, minimal feedback, and LLM-expert feedback across HumanEval, CoderEval, and SWE-bench. However, FeedbackEval excludes semantic static analysis (pylint/mypy)---it uses compiler output (syntax errors) rather than pylint-style warnings---and does not compare against execution test feedback with compute normalization. Our study directly addresses this gap: we add pylint/mypy as a treatment condition and hold token budget constant across all conditions.

\subsection{Type-Constrained Decoding}

\citet{Mundler2025Type} showed that type-constrained decoding---enforcing type correctness via vocabulary-filtered generation---reduces compilation errors by over 50\% on HumanEval and MBPP. This approach is fundamentally different from repair-mode feedback: it prevents certain errors at generation time rather than correcting them after. As a \emph{prevention} strategy, it is not directly comparable to iterative repair.

\subsection{Positioning Our Work}

Our study differs from prior work in three ways: (1) head-to-head pylint/mypy vs.\ execution comparison on the same benchmarks; (2) iso-compute control ($B$=1000 output tokens per problem across all conditions); (3) mechanism analysis decomposing pylint coverage by flag category.
